\documentclass[11pt,a4paper]{article}
\usepackage[utf8]{inputenc}
\usepackage[T1]{fontenc}
\usepackage{lmodern}
\usepackage[a4paper,margin=2.4cm]{geometry}
\usepackage{setspace,microtype,graphicx,booktabs,tabularx,array,longtable,enumitem,hyperref,xcolor,float,caption}
\usepackage{tikz}
\usepackage[backend=biber,style=authoryear,sorting=nyt,maxcitenames=2,maxbibnames=99]{biblatex}
\addbibresource{references.bib}
\usetikzlibrary{arrows.meta,positioning,shapes.geometric,calc}
\hypersetup{colorlinks=true,linkcolor=blue!50!black,citecolor=blue!50!black,urlcolor=blue!60!black,pdftitle={Responsible AI Technical Report: BrainWaveNet},pdfauthor={Levi Pedro Tejedor Rico}}
\setstretch{1.12}
\setlist[itemize]{leftmargin=*,itemsep=2pt,topsep=2pt}
\setlist[enumerate]{leftmargin=*,itemsep=2pt,topsep=2pt}
\newcolumntype{Y}{>{\raggedright\arraybackslash}X}
\newcommand{\ASD}{\textsc{asd}}
\newcommand{\TD}{\textsc{td}}
\newcommand{\PBG}{\textsc{pbg}}
\newcommand{\SUM}{\textsc{sum}}
\newcommand{\FAST}{\textsc{fast}}

\title{\textbf{Technical Report: Responsible AI Assessment of BrainWaveNet}\\
\large Responsible AI Assessment for Autism Spectrum Disorder Research}
\author{Levi Pedro Tejedor Rico\\Student ID: 4506001}
\date{Module: CSI\_7\_DEL -- Deep Learning}

\begin{document}
\maketitle

\begin{abstract}
This report assesses BrainWaveNet, a wavelet-based Transformer that classifies autism spectrum disorder (ASD) and typically developing participants using resting-state functional MRI. It combines Leslie's SUM values (Respect, Connect, Protect and Care), FAST principles (Fairness, Accountability, Sustainability and Transparency), and process-based governance across the AI lifecycle. BrainWaveNet is a research prototype evaluated on ABIDE-I; its reported performance and lack of external clinical validation do not support diagnostic use. The assessment identifies evidence gaps in stakeholder participation, subgroup evaluation, test-set governance, reproducibility and operational oversight. It proposes proportionate governance gates, accountable roles and social impact assessment checkpoints. The recommendation is to restrict the model to research until independent validation, meaningful autistic community participation and appropriate clinical safety governance have been demonstrated.
\end{abstract}

\tableofcontents
\newpage

\section{Introduction}
This report evaluates whether BrainWaveNet is suitable for use in a human healthcare environment and what would be required for responsible redevelopment. The project is a deep-learning research model that analyses resting-state functional magnetic resonance imaging (rs-fMRI) to classify participants as having an ASD diagnosis or belonging to a typically developing control group. The ethical question is not simply whether the model can distinguish these labels, but whether its purpose, evidence, design and governance justify use that could affect autistic people.

The assessment applies Leslie's SUM values and FAST principles, then translates them into a process-based governance (PBG) framework across problem formulation, data acquisition, preprocessing, modelling, validation and deployment. It distinguishes documented features of the published research from proposed controls: absence of information in a paper is treated as an evidence gap, not proof that an activity never occurred. The intended recommendation is therefore proportionate to the evidence available.

\section{AI Project Background}
\subsection{Clinical and data context}
Autism is a heterogeneous neurodevelopmental condition. Clinical identification involves developmental history and behavioural assessment; a neuroimaging classifier should not be assumed to provide a diagnosis or describe an individual's support needs. Research on functional connectivity has reported group-level patterns while also showing substantial complexity and overlap between individuals, limiting direct translation into individual diagnosis \parencite{holiga2019}.

The Autism Brain Imaging Data Exchange (ABIDE-I) shares structural and resting-state MRI data with phenotypic information. Its initial release included 1,112 participants from 17 sites: 539 with ASD and 573 typically developing controls, aged approximately 7--64 years \parencite{dimartino2014,abidewebsite}. Multi-site data broaden research opportunities but introduce variation in scanners, acquisition protocols, recruitment and participant characteristics. Site effects can confound neuroimaging analyses and require explicit evaluation \parencite{abraham2017,abideharmonization2022}.

\subsection{BrainWaveNet method and reported evidence}
Jeong et al. (2024) propose BrainWaveNet to learn spatiotemporal information from rs-fMRI in the frequency domain. A complex Morlet continuous wavelet transform (CWT) is applied to blood-oxygen-level-dependent time series. Residual convolutional blocks, frequency token encoding and temporal Transformer encoders process real and imaginary wavelet coefficients; spatial Transformer attention and cross-attention then learn relationships between regions of interest (ROIs). The output is a binary ASD/typically developing classification \parencite{jeong2024}.

The authors report CPAC-preprocessed data, a 0.01--0.1 Hz band-pass filter, no global signal regression, the Craddock 200 atlas and 78 time points per participant. CWT features use five frequency sub-bands, with sampling frequency adjusted for site-specific repetition time. Following quality control, 897 participants remained: 414 ASD and 483 controls. The reported protocol uses a 7:2:1 training, validation and test split and five-fold cross-validation, selecting models using validation AUC \parencite{jeong2024}.

BrainWaveNet reports AUC 0.7388, accuracy 67.55\%, sensitivity 66.49\% and specificity 68.35\%. It achieved the highest AUC, accuracy and sensitivity among the comparison methods reported, but not the highest specificity. Sensitivity and specificity also show notable standard deviations (9.17 and 9.80 percentage points). These are retrospective research results, not evidence of clinical utility, calibrated individual risk or improved patient outcomes. The study does not establish prospective performance or external validation on an independently collected cohort \parencite{jeong2024}.

The public repository supports inspection but does not, by itself, establish reproducibility. Its displayed results contain discrepancies from some paper values that should be reconciled, and the implementation requires verification in a documented clean environment \parencite{brainwavenetcode}. The paper's ROI analysis is exploratory; selected regions should not be presented as causal biomarkers or clinically validated explanations.

\section{Responsible AI Assessment}
The assessment uses SUM to judge the project's relationship with affected people and FAST to examine practical system properties. The lifecycle and PBG sections below convert the resulting gaps into controls. To avoid repeating the same critique, each principle is assessed here once; subsequent sections specify where and how the controls should operate.

\subsection{SUM values}
Leslie (2019) describes Respect, Connect, Protect and Care as values that support responsible AI. Table~\ref{tab:sum} summarises the evidence and the principal implication for BrainWaveNet.

\begin{table}[H]
\centering
\caption{Assessment of BrainWaveNet against SUM values.}
\label{tab:sum}
\small
\begin{tabularx}{\textwidth}{p{2.1cm}Y Y}
\toprule
\textbf{Value} & \textbf{Evidence and assessment} & \textbf{Required response}\\
\midrule
Respect & The dataset includes 414 autistic participants, but the paper does not report autistic co-design or their influence over research aims, interpretation or acceptable uses. Data representation is not equivalent to participation. ABIDE-I describes anonymisation and local ethics oversight, but neuroimaging linked to phenotype remains sensitive. & Establish accessible, paid autistic participation in decisions; respect neurodiversity in language; follow data permissions, minimise access and prohibit re-identification.\\
Connect & ABIDE-I connects 17 research sites and the paper makes methods and code inspectable. However, research collaboration does not demonstrate engagement with autistic people, families, clinicians or support organisations. & Co-design the use case with affected communities and clinicians. Consultation must be able to change or stop the project, not merely endorse a settled design.\\
Protect & The model is not shown making autonomous clinical decisions, limiting immediate direct risk. If used diagnostically, errors could cause distress, stigma, unnecessary follow-up or delayed support. No clinical security, incident or post-release safety arrangements are evidenced. & Restrict to research; before any pilot, establish a safety case, privacy and security controls, human oversight, complaint routes and incident response.\\
Care & The research question may contribute to understanding connectivity, but no patient-relevant benefit, improved diagnostic experience or better outcomes are demonstrated. Binary classification risks reducing a heterogeneous population to two labels. & Define a benefit with autistic people and clinicians; evaluate outcomes that matter to them; communicate uncertainty and never make access to support conditional on a model prediction.\\
\bottomrule
\end{tabularx}
\end{table}

\subsection{FAST principles}
FAST examines fairness, accountability, sustainability and transparency as operational requirements \parencite{leslie2019}. Table~\ref{tab:fast} focuses on what the available evidence can and cannot establish.

\begin{table}[H]
\centering
\caption{Assessment against FAST principles.}
\label{tab:fast}
\small
\begin{tabularx}{\textwidth}{p{2.1cm}Y Y}
\toprule
\textbf{Principle} & \textbf{Evidence gap or risk} & \textbf{Required control}\\
\midrule
Fairness & Overall diagnostic counts are reported, but not performance by age, sex, site or clinical subgroup. The 897-person quality-controlled sample is not compared in detail with excluded participants. Site and scanner effects may be learned instead of clinically relevant patterns \parencite{abraham2017,mehrabi2021}. & Publish inclusion/exclusion analyses and dataset factsheet; use site-held-out and external testing; report subgroup sensitivity, specificity, predictive values and calibration with uncertainty.\\
Accountability & Authors, paper and code are identifiable, but no operational clinical owner, complaints process or release authority is documented. The repository evaluates test metrics during training, creating a test-set peeking risk; this alone does not prove leakage. & Assign named accountable roles; seal the final test set until decisions are locked; preserve an audit trail and define approval, pause, appeal and withdrawal authority.\\
Sustainability & Energy, hardware use, inference cost and maintenance are not reported. MRI requires expensive specialist infrastructure and is not equally accessible. & Measure compute and energy; assess whole-pathway cost and MRI accessibility; compare expected benefit with lower-resource alternatives.\\
Transparency & The paper describes its method and the code is public, but metric discrepancies, incomplete reproducibility information and exploratory ROI interpretation limit independent scrutiny. No model card is evidenced. & Reconcile results; publish versioned code, environment, splits, seeds, preprocessing and model card; explain limitations in accessible language \parencite{mitchell2019}.\\
\bottomrule
\end{tabularx}
\end{table}

\section{Process-Based Governance and Lifecycle Assessment}
PBG embeds ethical review into project decisions rather than treating ethics as a one-off approval. The framework below is a proposed governance plan for redevelopment, not a claim that BrainWaveNet already operates these controls. Each lifecycle stage has an accountable owner, required evidence and a decision gate. SUM and FAST apply throughout.

\subsection{Problem formulation and stakeholder governance}
The published task is binary classification of ASD and typically developing participants from rs-fMRI. The paper does not establish a specific clinical user, care pathway or patient benefit. Before further development, an initial Social Impact Assessment (SIA) at the Alpha stage should establish whether the problem is worth pursuing, who may benefit or be harmed, and whether an MRI-based tool addresses a need identified by autistic people and clinicians. It should define prohibited uses, including autonomous diagnosis, and record alternatives such as not building the system.

A governance group should include an accountable project lead, clinical lead, data/ML lead, data-protection and security specialist, independent ethics/fairness reviewer, and autistic community representatives. Representatives should be appropriately supported and have real influence over purpose, evaluation, communication and release. A clinical governance or release board should be able to refuse, pause or withdraw the project.

\subsection{Data acquisition and preprocessing}
ABIDE-I began with 1,112 participants across 17 sites; BrainWaveNet retained 897 after quality control. The retained diagnostic groups are relatively balanced, but this does not establish representativeness across age, sex, site, IQ or clinical presentation. The paper does not provide a detailed subgroup comparison of included and excluded participants. ABIDE phenotype fields could support further analysis, although completeness varies. Original research permissions should not be assumed to authorise a new clinical purpose.

At the data gate, the team should document provenance, permitted use, ethics context, access controls, phenotype completeness and data minimisation. It should publish a participant flow diagram and compare retained and excluded cases by diagnosis, site and available demographic variables. At the preprocessing gate, versioned scripts should record CPAC version, filtering, parcellation, cropping, CWT settings and quality-control decisions. Sensitivity analyses should examine whether cropping to 78 volumes, motion, site and exclusion thresholds change performance. Harmonisation must be evaluated rather than assumed to remove site effects.

\subsection{Modelling, testing and validation}
The reported AUC of 0.7388 and accuracy of 67.55\% are encouraging only as preliminary research results. Sensitivity of 66.49\% and specificity of 68.35\% imply substantial classification error, and their standard deviations indicate variability. The paper does not provide external clinical validation, subgroup metrics or evidence that the model improves decisions. Its ROI ranking is exploratory and should not be interpreted as a causal explanation.

A responsible validation gate should lock the final test set before model selection, report confidence intervals and calibration, and evaluate site-held-out and independent external cohorts. Pre-specified subgroup analyses should cover age, sex, site and clinically relevant characteristics where sample sizes permit; small groups should be reported as uncertain rather than silently omitted. The team should compare clinically relevant baselines and determine statistical as well as clinical significance. The public implementation's repeated test evaluation should be removed from routine training logs or isolated so it cannot influence choices. Research on algorithmic bias shows that unequal representation and measurement processes can create systematic disparities, so evaluation must examine the full pipeline rather than accuracy alone \parencite{suresh2021,mehrabi2021}.

\subsection{Deployment, monitoring and SIA reassessment}
The available evidence describes research experiments and code publication, not routine hospital deployment, prospective clinical evaluation, regulatory approval or real-world monitoring. This is an evidence gap, not evidence of a failed clinical service. The model should therefore remain research-only and must not be presented as a diagnostic product.

If a future controlled pilot is proposed, an updated Alpha-to-Beta SIA should record changes made during development, stakeholder feedback, residual risks and reasons for proceeding. Release should require independent validation, institutional and regulatory review as applicable, clinical safety assessment, user training, privacy/security controls, human oversight and a documented rollback plan. Clinicians must retain decision authority and be able to disregard the output. Patients and families should receive clear explanations of intended use and uncertainty.

At the Beta stage, periodic reassessment should review subgroup performance, drift, site changes, complaints, adverse events, automation bias and unintended consequences. A named owner should have authority to suspend or withdraw the system when pre-agreed safety or fairness criteria are breached. SIA complements, rather than replaces, technical audit, clinical governance and monitoring \parencite{leslie2019,who2021,gmlp2021}.

\subsection{Governance gates and evidence}
Table~\ref{tab:gates} makes the proposed PBG process auditable. A gate should not be passed merely because documentation exists: evidence must be reviewed and the decision recorded.

\begin{longtable}{p{2.5cm}p{6.1cm}p{5.0cm}}
\caption{Proposed PBG gates for BrainWaveNet redevelopment.}
\label{tab:gates}\\
\toprule
\textbf{Stage} & \textbf{Evidence required} & \textbf{Gate decision}\\
\midrule
Alpha: purpose & Intended-use statement; stakeholder map; initial SIA; benefit hypothesis; prohibited uses; initial risk register. & Proceed only if affected stakeholders support a justified research question and risks are proportionate.\\
Data & Provenance, permissions, ethics context, privacy review, dataset factsheet and inclusion/exclusion analysis. & Reject data without clear permitted use; document representativeness and missingness limitations.\\
Preprocessing & Versioned pipeline, quality-control log, site/motion analysis and sensitivity tests. & Require reproducibility and evidence that processing does not create avoidable subgroup or site bias.\\
Model and validation & Model card, locked test protocol, confidence intervals, calibration, subgroup and external validation. & No clinical pilot until independent evidence meets pre-agreed clinical and fairness thresholds.\\
Alpha-to-Beta: pilot release & Updated SIA, clinical safety case, human-factors review, training, privacy/security, monitoring and rollback plans. & Release only under approved, supervised conditions; otherwise remain research-only.\\
Beta: reassessment & Monitoring reports, complaints, adverse events, drift and stakeholder review. & Continue, modify, pause or withdraw based on real-world evidence and documented review.\\
\bottomrule
\end{longtable}

\section{Examples of Non-Compliance and Corrective Priorities}
The following are the principal gaps supported by the available paper and repository evidence. They are described as missing evidence or implementation risks, not as proof of intentional misconduct.

\begin{itemize}
\item \textbf{No documented autistic co-design (SUM Respect/Connect):} research aims and interpretations may not reflect community priorities. \textit{Priority:} establish accessible, paid participation before defining a new intended use.
\item \textbf{No reported subgroup performance (FAST Fairness):} unequal error rates across age, sex, site or clinical groups could remain hidden. \textit{Priority:} publish subgroup analyses and uncertainty, alongside inclusion/exclusion comparisons.
\item \textbf{Test-set peeking risk (FAST Accountability):} the public implementation evaluates test metrics during training. This is not proof of leakage, but risks influencing choices. \textit{Priority:} isolate a locked test set and document all model-selection decisions.
\item \textbf{Incomplete reproducibility (FAST Transparency):} repository metrics differ from some paper values and the environment and pipeline need verification. \textit{Priority:} reconcile reported results and release a tested, versioned reproduction package.
\item \textbf{No clinical utility or deployment evidence (SUM Protect/Care):} reported performance does not demonstrate patient benefit or safe diagnostic use. \textit{Priority:} prohibit diagnosis and require independent clinical validation and safety governance before any pilot.
\item \textbf{Unmeasured resource burden (FAST Sustainability):} computational and MRI pathway costs are not reported. \textit{Priority:} measure energy, compute, access and maintenance against a credible benefit.
\end{itemize}

\section{Conclusion}
BrainWaveNet is a technically interesting research prototype, but the reported ABIDE-I results do not justify clinical diagnosis or autonomous decision-making. Its principal limitations are the absence of documented autistic co-design, subgroup and external clinical validation, complete reproducibility evidence, operational accountability and real-world monitoring. These are evidence and governance gaps; they should not be overstated as proof that harm has already occurred.

A proportionate next step is research-only redevelopment under the proposed PBG framework. The project should pass explicit purpose, data, preprocessing, validation and release gates, with SUM and FAST applied throughout and SIA conducted at Alpha, Alpha-to-Beta and Beta. Clinical use should remain prohibited unless autistic stakeholders and clinicians support the use case and independent evidence demonstrates acceptable safety, fairness and benefit.

\printbibliography[heading=bibintoc,title={References}]
\end{document}
